% ============================================================
%  02_revision_literatura.tex — Capítulo 2: Revisión de literatura
% ============================================================

\chapter{Revisión de literatura y trabajos relacionados}
\label{cap:estado-del-arte}

\textit{[Este capítulo resume la literatura relacionada. Usa citas con \texttt{\textbackslash cite\{clave\}}
y añade las referencias a \texttt{bibliografia/referencias.bib}.]}

\section{Marco conceptual}
\label{sec:marco-conceptual}

\textit{[Definir los conceptos que sustentan la propuesta: cítricos, invernaderos,
\textit{Tetranychus urticae}, temperatura, humedad relativa, instrumentación e IoT.
Presentar las bases de ontologías y representación del conocimiento que sean
pertinentes para el proyecto.]}

\section{Estado del arte}
\label{sec:estado-del-arte}

\textit{[Describir la búsqueda y selección de literatura sobre monitoreo ambiental,
agricultura de precisión y detección de la plaga. Incorporar únicamente referencias
verificadas en la bibliografía.]}

\subsection{Trabajos relacionados}
\label{subsec:trabajos-relacionados}

\textit{[Comparar los trabajos seleccionados: objetivos, arquitecturas, sensores,
protocolos, métodos de análisis y resultados reportados.]}

\subsection{Brechas existentes}
\label{subsec:brechas}

\textit{[Identificar limitaciones y vacíos sustentados en la revisión de literatura.]}

\subsection{Aportes del proyecto}
\label{subsec:aportes-proyecto}

\textit{[Explicar los aportes propuestos frente a las brechas identificadas, sin
presentarlos como resultados ya demostrados.]}
